\documentclass[11pt]{article}
\usepackage{amsmath, amssymb, geometry, xcolor, tcolorbox}
\geometry{a4paper, margin=1in}

\title{\textbf{Homework 4.3: Auxiliary Equations with Complex Roots}}
\author{Tashfeen Omran}
\date{September 2026}

\begin{document}

\maketitle
\tableofcontents
\newpage

\section{Part 1: Comprehensive Introduction, Context, and Prerequisites}

\subsection{Core Concepts}
In the study of linear second-order homogeneous differential equations with constant coefficients, $ay'' + by' + cy = 0$, we rely on the auxiliary (or characteristic) equation $ar^2 + br + c = 0$. When the discriminant is negative ($b^2 - 4ac < 0$), the roots of this quadratic equation are complex conjugates: $r = \alpha \pm i\beta$, where $\alpha$ and $\beta$ are real numbers and $i = \sqrt{-1}$. 

This chapter explores how to interpret these complex roots to form real-valued solutions. The core concept is that a complex exponential solution $e^{(\alpha + i\beta)t}$ can be split into real and imaginary parts, both of which independently satisfy the original differential equation.

\subsection{Intuition and Derivation}
If we blindly apply the rule from real roots, we get a solution of the form $y(t) = C_1 e^{(\alpha + i\beta)t} + C_2 e^{(\alpha - i\beta)t}$. However, in physical and financial modeling, we require real-valued functions. 
To bridge this gap, we use \textbf{Euler's Formula}: $e^{i\theta} = \cos\theta + i\sin\theta$. 
Applying this to our complex exponential:
\[ e^{(\alpha + i\beta)t} = e^{\alpha t} e^{i\beta t} = e^{\alpha t} (\cos \beta t + i\sin \beta t) \]
Because the differential equation is linear and homogeneous, the real part $u(t) = e^{\alpha t} \cos \beta t$ and the imaginary part $v(t) = e^{\alpha t} \sin \beta t$ are both valid, linearly independent solutions. Thus, the general real-valued solution is their linear combination.

\subsection{Historical Context and Motivation}
The foundational work for this topic was laid by Leonhard Euler in his 1748 masterpiece \textit{Introductio in analysin infinitorum}. Mathematicians of the 18th century were grappling with the mechanics of vibrating strings and pendulums. They knew these systems oscillated, but the algebraic methods of the time yielded "imaginary" exponential solutions. Euler's brilliant insight to connect complex exponentials with trigonometric functions not only solved these differential equations but also unified algebra and trigonometry, forever changing mathematical analysis.

\subsection{Key Formulas}
\begin{itemize}
    \item \textbf{Auxiliary Equation:} $ar^2 + br + c = 0$
    \item \textbf{Complex Roots:} $r = \alpha \pm i\beta = -\frac{b}{2a} \pm i\frac{\sqrt{4ac - b^2}}{2a}$
    \item \textbf{Euler's Formula:} $e^{i\theta} = \cos\theta + i\sin\theta$
    \item \textbf{General Solution:} $y(t) = e^{\alpha t}(c_1 \cos \beta t + c_2 \sin \beta t)$
\end{itemize}

\subsection{Prerequisites}
\begin{itemize}
    \item \textbf{Algebra:} Proficiency with the quadratic formula and complex number arithmetic.
    \item \textbf{Calculus:} The Product Rule and Chain Rule for differentiation (crucial for Initial Value Problems).
    \item \textbf{Pre-Calculus:} Basic trigonometric identities and properties of sine and cosine.
\end{itemize}

\newpage
\section{Part 2: Detailed Homework Solutions}
\textit{Strategic Selection Rationale: To maximize learning efficiency, I have selected a representative sample of problems from the provided text that covers every distinct archetype: a pure imaginary root problem (Simple Harmonic), a complex root problem (Damped Harmonic), an Initial Value Problem (IVP), a higher-order equation, and a real-world modeling problem (RLC Circuit). This ensures full conceptual coverage without redundant repetition.}

\subsection{Archetype 1: Pure Imaginary Roots (Problem 1)}
\textbf{Problem:} Find a general solution to $y'' + 9y = 0$.

\textbf{Solution:}
\begin{enumerate}
    \item \textbf{Set up the auxiliary equation:} 
    \[ r^2 + 9 = 0 \]
    \item \textbf{Solve for $r$:}
    \[ r^2 = -9 \implies r = \pm 3i \]
    Here, the real part $\alpha = 0$ and the imaginary part $\beta = 3$.
    \item \textbf{Formulate the general solution:}
    Using the formula $y(t) = e^{\alpha t}(c_1 \cos \beta t + c_2 \sin \beta t)$:
    \[ y(t) = e^{0 \cdot t}(c_1 \cos 3t + c_2 \sin 3t) \]
    \[ y(t) = c_1 \cos 3t + c_2 \sin 3t \]
\end{enumerate}

\subsection{Archetype 2: Complex Roots with Real Part (Problem 13)}
\textbf{Problem:} Find a general solution to $y'' - 2y' + 26y = 0$.

\textbf{Solution:}
\begin{enumerate}
    \item \textbf{Set up the auxiliary equation:} 
    \[ r^2 - 2r + 26 = 0 \]
    \item \textbf{Solve for $r$ using the quadratic formula:}
    \[ r = \frac{-(-2) \pm \sqrt{(-2)^2 - 4(1)(26)}}{2(1)} = \frac{2 \pm \sqrt{4 - 104}}{2} \]
    \[ r = \frac{2 \pm \sqrt{-100}}{2} = \frac{2 \pm 10i}{2} = 1 \pm 5i \]
    Here, $\alpha = 1$ and $\beta = 5$.
    \item \textbf{Formulate the general solution:}
    \[ y(t) = e^{t}(c_1 \cos 5t + c_2 \sin 5t) \]
\end{enumerate}

\subsection{Archetype 3: Initial Value Problem (Problem 21)}
\textbf{Problem:} Solve the initial value problem $y'' + 2y' + 2y = 0; \quad y(0) = 2, \quad y'(0) = 1$.

\textbf{Solution:}
\begin{enumerate}
    \item \textbf{Find the general solution:}
    \[ r^2 + 2r + 2 = 0 \implies r = \frac{-2 \pm \sqrt{4 - 8}}{2} = \frac{-2 \pm 2i}{2} = -1 \pm i \]
    Here, $\alpha = -1$ and $\beta = 1$.
    \[ y(t) = e^{-t}(c_1 \cos t + c_2 \sin t) \]
    \item \textbf{Apply the first initial condition $y(0) = 2$:}
    \[ y(0) = e^{0}(c_1 \cos 0 + c_2 \sin 0) = 1 \cdot (c_1 \cdot 1 + c_2 \cdot 0) = c_1 \]
    \[ c_1 = 2 \]
    So, $y(t) = e^{-t}(2 \cos t + c_2 \sin t)$.
    \item \textbf{Differentiate $y(t)$ using the Product Rule:}
    \[ y'(t) = \frac{d}{dt}[e^{-t}] \cdot (2 \cos t + c_2 \sin t) + e^{-t} \cdot \frac{d}{dt}[2 \cos t + c_2 \sin t] \]
    \[ y'(t) = -e^{-t}(2 \cos t + c_2 \sin t) + e^{-t}(-2 \sin t + c_2 \cos t) \]
    \item \textbf{Apply the second initial condition $y'(0) = 1$:}
    \[ y'(0) = -1(2 \cos 0 + c_2 \sin 0) + 1(-2 \sin 0 + c_2 \cos 0) \]
    \[ 1 = -1(2) + 1(c_2) \implies 1 = -2 + c_2 \implies c_2 = 3 \]
    \item \textbf{State the final solution:}
    \[ y(t) = e^{-t}(2 \cos t + 3 \sin t) \]
\end{enumerate}

\subsection{Archetype 4: Higher-Order Equations (Problem 29a)}
\textbf{Problem:} Find a general solution to $y''' - y'' + y' + 3y = 0$.

\textbf{Solution:}
\begin{enumerate}
    \item \textbf{Set up the auxiliary equation:}
    \[ r^3 - r^2 + r + 3 = 0 \]
    \item \textbf{Find the roots:}
    By the Rational Root Theorem, we test factors of 3. Let's test $r = -1$:
    \[ (-1)^3 - (-1)^2 + (-1) + 3 = -1 - 1 - 1 + 3 = 0 \]
    Since $r = -1$ is a root, $(r+1)$ is a factor. We use polynomial long division or synthetic division to divide $r^3 - r^2 + r + 3$ by $(r+1)$, yielding $r^2 - 2r + 3$.
    Now, solve $r^2 - 2r + 3 = 0$:
    \[ r = \frac{2 \pm \sqrt{4 - 12}}{2} = \frac{2 \pm \sqrt{-8}}{2} = \frac{2 \pm 2i\sqrt{2}}{2} = 1 \pm i\sqrt{2} \]
    \item \textbf{Formulate the general solution:}
    The real root $r = -1$ gives $c_1 e^{-t}$. The complex roots $1 \pm i\sqrt{2}$ give $e^t(c_2 \cos(\sqrt{2}t) + c_3 \sin(\sqrt{2}t))$.
    \[ y(t) = c_1 e^{-t} + e^t(c_2 \cos(\sqrt{2}t) + c_3 \sin(\sqrt{2}t)) \]
\end{enumerate}

\subsection{Archetype 5: Physical Application - RLC Circuit (Problem 34)}
\textbf{Problem:} An RLC circuit is governed by $L \frac{dI}{dt} + RI + \frac{q}{C} = E(t)$. Differentiate to get $L I'' + R I' + \frac{1}{C} I = E'(t)$. Given $L = 10\text{ H}$, $R = 20\text{ }\Omega$, $C = (6260)^{-1}\text{ F}$, and $E(t) = 100\text{ V}$. Find the current $I(t)$ if $q(0) = 0$ and $I(0) = 0$.

\textbf{Solution:}
\begin{enumerate}
    \item \textbf{Set up the differential equation:}
    Since $E(t) = 100$ is constant, $E'(t) = 0$. Substituting the constants:
    \[ 10 I'' + 20 I' + 6260 I = 0 \implies I'' + 2I' + 626I = 0 \]
    \item \textbf{Find the general solution:}
    \[ r^2 + 2r + 626 = 0 \implies r = \frac{-2 \pm \sqrt{4 - 2504}}{2} = \frac{-2 \pm \sqrt{-2500}}{2} = -1 \pm 25i \]
    \[ I(t) = e^{-t}(c_1 \cos 25t + c_2 \sin 25t) \]
    \item \textbf{Apply initial conditions:}
    We are given $I(0) = 0$.
    \[ I(0) = e^0(c_1 \cos 0 + c_2 \sin 0) = c_1 \implies c_1 = 0 \]
    So, $I(t) = c_2 e^{-t} \sin 25t$.
    
    To find $I'(0)$, we use the original equation at $t=0$: $L I'(0) + R I(0) + \frac{q(0)}{C} = E(0)$.
    \[ 10 I'(0) + 20(0) + 6260(0) = 100 \implies 10 I'(0) = 100 \implies I'(0) = 10 \]
    
    Now, differentiate $I(t) = c_2 e^{-t} \sin 25t$:
    \[ I'(t) = c_2 \left( -e^{-t} \sin 25t + 25e^{-t} \cos 25t \right) \]
    \[ I'(0) = c_2 ( -0 + 25(1) ) = 25c_2 \]
    Set $I'(0) = 10$:
    \[ 25c_2 = 10 \implies c_2 = \frac{10}{25} = 0.4 \]
    \item \textbf{State the final solution:}
    \[ I(t) = 0.4 e^{-t} \sin 25t \]
\end{enumerate}

\newpage
\section{Part 3: In-Depth Analysis of Problems and Techniques}

\subsection{A) Problem Types and General Approach}
\begin{itemize}
    \item \textbf{Pure Imaginary Roots (Simple Harmonic):} The auxiliary equation has no $r$ term (e.g., $r^2 + k^2 = 0$). The strategy is to recognize $\alpha = 0$, meaning there is no exponential decay/growth. The solution is purely oscillatory: $y(t) = c_1 \cos \beta t + c_2 \sin \beta t$.
    \item \textbf{Complex Roots with Real Part (Damped Harmonic):} The auxiliary equation yields $r = \alpha \pm i\beta$. The strategy is to identify $\alpha$ (the damping factor) and $\beta$ (the frequency factor) and plug them directly into the general formula.
    \item \textbf{Initial Value Problems (IVPs):} These require finding the specific constants $c_1$ and $c_2$. The strategy is to first find the general solution, apply $y(0)$ to find $c_1$, take the derivative using the Product Rule, and apply $y'(0)$ to find $c_2$.
    \item \textbf{Higher-Order Equations:} The strategy is to factor the characteristic polynomial (often using the Rational Root Theorem and synthetic division) to break it down into linear and irreducible quadratic factors, then apply the rules for real and complex roots respectively.
\end{itemize}

\subsection{B) Key Algebraic and Calculus Manipulations}
\begin{itemize}
    \item \textbf{The Product Rule for IVPs:} In Problem 21, differentiating $y(t) = e^{-t}(c_1 \cos t + c_2 \sin t)$ requires the Product Rule. A common mistake is to only differentiate the inside of the parentheses. The correct manipulation is $\frac{d}{dt}[uv] = u'v + uv'$.
    \item \textbf{Extracting $I'(0)$ from the Physical Equation:} In Problem 34, we were not explicitly given $I'(0)$. The crucial algebraic trick was to use the un-differentiated equation $L I'(t) + R I(t) + \frac{q(t)}{C} = E(t)$ evaluated at $t=0$ to algebraically solve for $I'(0)$.
    \item \textbf{Rational Root Theorem:} In Problem 29a, solving a cubic equation $r^3 - r^2 + r + 3 = 0$ required guessing a root by checking factors of the constant term (3). Finding $r=-1$ allowed us to factor the cubic into a linear and a quadratic term.
\end{itemize}

\section{Part 4: "Cheatsheet" and Tips for Success}

\subsection{Summary of Important Formulas}
\begin{tcolorbox}
\textbf{Roots:} $r = \alpha \pm i\beta$ \\
\textbf{General Solution:} $y(t) = e^{\alpha t}(c_1 \cos \beta t + c_2 \sin \beta t)$ \\
\textbf{Derivative of General Solution:} \\
$y'(t) = e^{\alpha t} [ (\alpha c_1 + \beta c_2)\cos \beta t + (\alpha c_2 - \beta c_1)\sin \beta t ]$
\end{tcolorbox}

\subsection{Tricks and Shortcuts}
\begin{itemize}
    \item \textbf{The $t=0$ Shortcut:} When evaluating initial conditions at $t=0$, the sine terms vanish and cosine terms become 1. Thus, $y(0) = c_1$. 
    \item \textbf{The $y'(0)$ Shortcut:} If you memorize the derivative structure, you can instantly know that $y'(0) = \alpha c_1 + \beta c_2$. This saves you from doing the Product Rule on the exam!
\end{itemize}

\subsection{Common Pitfalls}
\begin{itemize}
    \item \textbf{Forgetting the $i$:} When $b^2 - 4ac$ is negative, students often pull out the negative sign but forget to write $i$, leading to incorrect real roots.
    \item \textbf{Confusing $\alpha$ and $\beta$:} $\alpha$ is the real part (goes in the exponent), $\beta$ is the imaginary part (goes inside the trig functions). Do not swap them!
    \item \textbf{Including $i$ in the final answer:} The general solution $y(t) = e^{\alpha t}(c_1 \cos \beta t + c_2 \sin \beta t)$ is strictly real. Never write $\cos(i\beta t)$ or include $i$ in the constants unless specifically dealing with complex vector spaces.
\end{itemize}

\newpage
\section{Part 5: Conceptual Synthesis and The "Big Picture"}

\subsection{A) Thematic Connections}
The central theme of this topic is \textbf{Linearity and Superposition}. Just as any vector in a 2D plane can be described as a combination of $\hat{i}$ and $\hat{j}$, any solution to a second-order linear homogeneous ODE can be described as a linear combination of two fundamental, linearly independent solutions. Euler's formula allows us to take complex basis vectors ($e^{(\alpha + i\beta)t}$) and transform them into a real basis ($e^{\alpha t}\cos \beta t$ and $e^{\alpha t}\sin \beta t$) that spans the exact same solution space.

\subsection{B) Forward and Backward Links}
\begin{itemize}
    \item \textbf{Backward Links:} This topic is the logical completion of the previous chapter on real roots. It relies heavily on the algebraic foundations of the quadratic formula and the calculus of exponential functions.
    \item \textbf{Forward Links:} Mastering these oscillatory solutions is the absolute bedrock for upcoming topics. When we introduce \textit{Forced Oscillations} (non-homogeneous equations), these homogeneous solutions will represent the "transient" state of the system. Furthermore, understanding how exponentials and sinusoids interact is the fundamental prerequisite for \textbf{Laplace Transforms} and \textbf{Fourier Series}.
\end{itemize}

\section{Part 6: Real-World Application and Modeling}

\subsection{A) Concrete Scenarios in Finance and Economics}
\begin{enumerate}
    \item \textbf{Macroeconomic Business Cycles (Multiplier-Accelerator Model):} In economics, the interaction between the Keynesian multiplier (consumer spending) and the investment accelerator (capital investment) can be modeled using a second-order ODE. If the roots of the characteristic equation are complex, the model predicts endogenous business cycles—periods of economic boom and bust that naturally oscillate and dampen over time.
    \item \textbf{Stochastic Calculus and Mean-Reverting Asset Prices:} In quantitative finance, certain commodities (like natural gas or agricultural products) exhibit seasonality and mean-reversion. While typically modeled with stochastic differential equations (SDEs), the deterministic drift component often takes the form of a damped harmonic oscillator to model how prices "spring" back to an equilibrium level after a supply shock, overshooting and oscillating before settling.
\end{enumerate}

\subsection{B) Model Problem Setup: The Business Cycle}
Let $Y(t)$ represent the deviation of a nation's GDP from its long-term equilibrium trend. Economic theorists propose that the rate of change of GDP acceleration is negatively proportional to both its current rate of change (due to friction/inflationary pressures, with coefficient $\delta$) and its current level (due to market saturation, with coefficient $\omega^2$).

\textbf{Variables:}
\begin{itemize}
    \item $Y(t)$: GDP deviation from equilibrium at time $t$.
    \item $\delta$: The economic damping coefficient (friction).
    \item $\omega^2$: The economic stiffness (restoring force of the market).
\end{itemize}

\textbf{Mathematical Model:}
The economy is governed by the differential equation:
\[ Y'' + \delta Y' + \omega^2 Y = 0 \]
To find the trajectory of the GDP after a sudden market crash where $Y(0) = -Y_0$ and $Y'(0) = 0$, an economist would solve the auxiliary equation $r^2 + \delta r + \omega^2 = 0$. If $\delta^2 < 4\omega^2$, the roots are complex, and the economy will experience damped oscillations (boom and bust cycles) as it returns to equilibrium.

\newpage
\section{Part 7: Common Variations and Untested Concepts}

\subsection{Boundary Value Problems (BVPs)}
\textbf{Concept:} The provided homework exclusively featured Initial Value Problems (IVPs), where conditions are given at a single point (e.g., $t=0$). A common variation is a Boundary Value Problem (BVP), where conditions are specified at \textit{two different points} (e.g., $t=0$ and $t=L$). BVPs are crucial in partial differential equations and quantum mechanics. Unlike IVPs, BVPs might have unique solutions, infinite solutions, or no solution at all.

\textbf{Worked Example:}
Solve $y'' + 4y = 0$ subject to $y(0) = 0$ and $y(\pi) = 0$.

\textbf{Solution:}
1. The general solution is $y(t) = c_1 \cos 2t + c_2 \sin 2t$. \\
2. Apply the first boundary condition $y(0) = 0$:
\[ y(0) = c_1 \cos(0) + c_2 \sin(0) = c_1 \implies c_1 = 0 \]
So, $y(t) = c_2 \sin 2t$. \\
3. Apply the second boundary condition $y(\pi) = 0$:
\[ y(\pi) = c_2 \sin(2\pi) = c_2(0) = 0 \]
Notice that $0 = 0$ is true for \textit{any} value of $c_2$. \\
\textbf{Conclusion:} This BVP has infinitely many solutions of the form $y(t) = c_2 \sin 2t$, where $c_2$ is any real number.

\section{Part 8: Advanced Diagnostic Testing: "Find the Flaw"}
\textit{Instructions: Act as the grader. For each problem, identify the specific error, explain why it is an error in one sentence, and provide the correct final solution.}

\subsection*{Problem 1}
\textbf{Task:} Solve $y'' + 16y = 0$. \\
\textbf{Flawed Solution:} 
The auxiliary equation is $r^2 + 16 = 0$. 
Solving for $r$, we get $r^2 = 16 \implies r = \pm 4$. 
Therefore, the general solution is $y(t) = c_1 e^{4t} + c_2 e^{-4t}$.

\subsection*{Problem 2}
\textbf{Task:} Solve $y'' - 6y' + 25y = 0$. \\
\textbf{Flawed Solution:} 
The auxiliary equation is $r^2 - 6r + 25 = 0$. 
Using the quadratic formula: $r = \frac{6 \pm \sqrt{36 - 100}}{2} = \frac{6 \pm \sqrt{-64}}{2} = 3 \pm 4i$. 
Here, $\alpha = 4$ and $\beta = 3$. 
The general solution is $y(t) = e^{4t}(c_1 \cos 3t + c_2 \sin 3t)$.

\subsection*{Problem 3}
\textbf{Task:} Solve the IVP $y'' + 2y' + 5y = 0, \quad y(0)=1, \quad y'(0)=0$. \\
\textbf{Flawed Solution:} 
$r^2 + 2r + 5 = 0 \implies r = -1 \pm 2i$. 
General solution: $y(t) = e^{-t}(c_1 \cos 2t + c_2 \sin 2t)$. 
Using $y(0)=1$: $c_1 = 1$. 
Now, take the derivative: $y'(t) = e^{-t}(-2c_1 \sin 2t + 2c_2 \cos 2t)$. 
Using $y'(0)=0$: $e^0(-0 + 2c_2) = 0 \implies 2c_2 = 0 \implies c_2 = 0$. 
Final answer: $y(t) = e^{-t} \cos 2t$.

\subsection*{Problem 4}
\textbf{Task:} Solve the IVP $y'' + 9y = 0, \quad y(0)=2, \quad y'(0)=3$. \\
\textbf{Flawed Solution:} 
$r^2 + 9 = 0 \implies r = \pm 3i$. 
General solution: $y(t) = c_1 \cos 3t + c_2 \sin 3t$. 
Using $y(0)=2$: $c_1 = 2$. 
Derivative: $y'(t) = -3c_1 \sin 3t + c_2 \cos 3t$. 
Using $y'(0)=3$: $c_2 = 3$. 
Final answer: $y(t) = 2 \cos 3t + 3 \sin 3t$.

\subsection*{Problem 5}
\textbf{Task:} Find the general solution to $y''' + y' = 0$. \\
\textbf{Flawed Solution:} 
The auxiliary equation is $r^3 + 1 = 0$. 
One obvious root is $r = -1$. 
Factoring gives $(r+1)(r^2 - r + 1) = 0$. 
The quadratic roots are $r = \frac{1 \pm \sqrt{1 - 4}}{2} = \frac{1}{2} \pm i\frac{\sqrt{3}}{2}$. 
Final answer: $y(t) = c_1 e^{-t} + e^{t/2}\left(c_2 \cos\left(\frac{\sqrt{3}}{2}t\right) + c_3 \sin\left(\frac{\sqrt{3}}{2}t\right)\right)$.

\vspace{1cm}
\hrule
\vspace{0.5cm}
\textbf{Grader's Answer Key (For your reference):}
\begin{itemize}
    \item \textbf{Problem 1 Flaw:} The student missed the negative sign when moving 16 to the other side; $r^2 = -16$, so the roots are imaginary ($\pm 4i$), not real. \textbf{Correct:} $y(t) = c_1 \cos 4t + c_2 \sin 4t$.
    \item \textbf{Problem 2 Flaw:} The student swapped the real part $\alpha$ and the imaginary part $\beta$. \textbf{Correct:} $y(t) = e^{3t}(c_1 \cos 4t + c_2 \sin 4t)$.
    \item \textbf{Problem 3 Flaw:} The student forgot to use the Product Rule when differentiating $y(t)$. \textbf{Correct:} $y'(t) = -e^{-t}(\cos 2t + c_2 \sin 2t) + e^{-t}(-2 \sin 2t + 2c_2 \cos 2t)$. Setting $y'(0)=0$ gives $-1 + 2c_2 = 0 \implies c_2 = 1/2$. Final: $y(t) = e^{-t}(\cos 2t + \frac{1}{2}\sin 2t)$.
    \item \textbf{Problem 4 Flaw:} The student forgot the Chain Rule on the sine term when differentiating; the derivative of $c_2 \sin 3t$ is $3c_2 \cos 3t$, not $c_2 \cos 3t$. \textbf{Correct:} $3c_2 = 3 \implies c_2 = 1$. Final: $y(t) = 2 \cos 3t + \sin 3t$.
    \item \textbf{Problem 5 Flaw:} The student wrote the wrong auxiliary equation; $y'$ corresponds to $r$, not $1$, so the equation should be $r^3 + r = 0$. \textbf{Correct:} $r(r^2+1)=0 \implies r=0, \pm i$. Final: $y(t) = c_1 + c_2 \cos t + c_3 \sin t$.
\end{itemize}

\end{document}